% Copyright 2026 Myeongjin Kang. Original prose CC BY 4.0; original code MIT.
\documentclass[11pt,a4paper]{article}
\usepackage[T1]{fontenc}
\usepackage{lmodern}
\usepackage[margin=1in]{geometry}
\usepackage{amsmath,amssymb,amsthm,mathtools}
\usepackage{microtype,booktabs,array}
\usepackage{xcolor}
\usepackage[colorlinks=true,linkcolor=blue!50!black,citecolor=blue!50!black,urlcolor=blue!50!black]{hyperref}
\hypersetup{pdftitle={The Borsuk number of the complex projective plane is nine},pdfauthor={Myeongjin Kang},pdfsubject={Complex projective coverings and rank-one Hermitian projectors}}
\setlength{\emergencystretch}{2em}
\newtheorem{maintheorem}{Theorem}
\renewcommand{\themaintheorem}{\Alph{maintheorem}}
\newtheorem{theorem}{Theorem}[section]
\newtheorem{lemma}[theorem]{Lemma}
\newtheorem{proposition}[theorem]{Proposition}
\newtheorem{corollary}[theorem]{Corollary}
\theoremstyle{remark}
\newtheorem{remark}[theorem]{Remark}
\numberwithin{equation}{section}
\newcommand{\C}{\mathbb C}
\newcommand{\R}{\mathbb R}
\newcommand{\CP}{\mathbb{CP}}
\newcommand{\Herm}{\operatorname{Herm}}
\newcommand{\tr}{\operatorname{tr}}
\newcommand{\ran}{\operatorname{ran}}
\newcommand{\diam}{\operatorname{diam}}
\newcommand{\ip}[2]{\langle #1,#2\rangle}
\newcommand{\norm}[1]{\lVert #1\rVert}
\newcommand{\abs}[1]{\lvert #1\rvert}
\title{The Borsuk number of the complex projective plane is nine}
\author{Myeongjin Kang}
\date{October 2026}
\begin{document}
\maketitle
\begin{abstract}
Let $X=\{uu^*:u\in\C^3,\ \norm{u}=1\}$ be the rank-one projector realization of the complex projective plane in the eight-dimensional affine space of trace-one Hermitian matrices, with the Frobenius metric. We prove that $\diam X=\sqrt2$ and that its Borsuk number is exactly nine. The lower bound applies to covers by arbitrary subsets and follows from a Hermitian version of a known spectral odd-map construction. For the upper bound, we perturb the nine rays of the Hesse symmetric informationally complete configuration. A cyclic assignment of its twelve mutually unbiased basis poles is compatible with the Weyl--Heisenberg action. A single covariant tangent, together with a uniform two-trio estimate in full complex neighborhoods, removes all equality holes and excludes displaced ones. The resulting nine maximum-overlap cells are compact and have strictly smaller diameter. The same construction gives nine Fubini--Study balls of radius strictly below $\pi/4$ covering $\CP^2$. The result attains the bound $d+1$ in dimension eight.
\end{abstract}

\section{Statement and relation to earlier methods}
For a bounded metric space $Y$ of positive diameter, its \emph{Borsuk number} $b(Y)$ is the least number of subsets of $Y$, each of diameter strictly smaller than $\diam Y$, whose union is $Y$. No regularity is imposed on the subsets. We use the Hermitian inner product $\ip uv=u^*v$, conjugate-linear in its first argument, and put
\[
 \mathcal T=\{A\in\Herm_3:\tr A=1\},\qquad
 X=\{P_u=uu^*:u\in\C^3,\ \norm u=1\}.
\]
The real vector space $\Herm_3$ has dimension nine, and $\mathcal T$ has dimension eight. Its metric is induced by the Frobenius norm $\norm A_F^2=\tr(A^*A)$.

\begin{maintheorem}\label{thm:main}
The space $X$ has diameter $\sqrt2$ and Borsuk number nine. More precisely:
\begin{enumerate}
\item There exist nine compact subsets covering $X$, each of diameter strictly smaller than $\sqrt2$.
\item In every cover of $X$ by at most eight arbitrary subsets, some member has diameter $\sqrt2$.
\end{enumerate}
\end{maintheorem}
Thus this projector manifold attains $b(X)=9=\dim(\mathcal T)+1$; it is not a counterexample to Borsuk's covering assertion in dimension eight.

For a finite collection $C$ of unit rays define
\begin{equation}\label{eq:alpha}
 \alpha(C)=\min_{[u]\in\CP^2}\max_{s\in C}\abs{\ip su}^2.
\end{equation}
Our upper construction gives nine rays with $\alpha(C)>1/2$. With the convention $d_{\mathrm{FS}}([u],[v])=\arccos\abs{\ip uv}$, it therefore gives nine closed balls of radius strictly less than $\pi/4$ covering $\CP^2$. The distinction between a strict uniform bound in \eqref{eq:alpha} and a pairwise packing bound is essential.

The lower-bound mechanism is adapted from the real-projector construction of OpenAI~\cite{OpenAI}: enlarge a strict-gap cover, extend a homogeneous partition to positive semidefinite matrices, and subtract the extensions on the positive and negative spectral parts. We give the complete Hermitian argument. No real equality-case obstruction or source theorem is needed as a premise. The Hesse configuration, Weyl--Heisenberg action, and twelve mutually unbiased normals are standard antecedents; see Dang, Blanchfield, Bengtsson and Appleby~\cite{Dang} and Hughston and Salamon~\cite{HS}. Their geometry is recalled explicitly below.

The additional operation is the uniform removal of the Hesse equality holes by a non-SIC covariant perturbation. A positive derivative at each old hole alone would not control new nearby minimizers. Lemma~\ref{lem:local} supplies that missing neighborhood estimate. Compactness and the conversion to strict smaller-diameter cells are then standard. Section~\ref{sec:comparison} and Appendix~\ref{app:nonsic} distinguish this operation from specific earlier constructions. This is a bounded comparison, not a claim to have classified all nine-ray coverings or established worldwide priority.

\section{Projector geometry and the arbitrary-cover lower bound}
\begin{lemma}\label{lem:geometry}
For unit vectors $u,v\in\C^3$,
\begin{equation}\label{eq:projector}
 \norm{P_u-P_v}_F^2=2-2\abs{\ip uv}^2.
\end{equation}
Consequently $\diam X=\sqrt2$, with equality in the pairwise distance exactly for orthogonal rays. The map $[u]\mapsto P_u$ is a homeomorphism $\CP^2\to X$.
\end{lemma}
\begin{proof}
We have $\tr(P_u^2)=1$ and $\tr(P_uP_v)=\abs{\ip uv}^2$, giving \eqref{eq:projector}. Orthogonal unit vectors exist. The projector map is continuous and injective, since its range determines the line. It is a homeomorphism from compact $\CP^2$ onto its image in a Hausdorff space.
\end{proof}

\begin{lemma}[Strict-gap enlargement]\label{lem:gap}
Suppose $X=A_1\cup\cdots\cup A_m$ and $\diam A_i<\sqrt2$ for every nonempty member. There are continuous functions $f_i:\CP^2\to[0,1]$ with $\sum_i f_i=1$ such that $\{f_i>0\}$ contains no pair of orthogonal rays.
\end{lemma}
\begin{proof}
For nonempty $A_i$ set $t_i=(\sqrt2-\diam A_i)/4$ and let $U_i$ be its relative open $t_i$-neighborhood in $X$. For empty $A_i$ put $U_i=\varnothing$. Approximation by points of $A_i$ and the triangle inequality give
\[
 \diam U_i\le\diam A_i+2t_i<\sqrt2.
\]
Thus each $U_i$ is proper, its complement is nonempty and closed, and it contains no orthogonal pair. Let $g_i(P)=\operatorname{dist}(P,X\setminus U_i)$. These continuous functions are positive exactly on $U_i$, and their sum is everywhere positive because the $U_i$ cover $X$. Normalize by $f_i=g_i/\sum_jg_j$ and use Lemma~\ref{lem:geometry}. This argument imposes no closure or measurability requirement on $A_i$.
\end{proof}
For nonzero $w$, write $S(w)=\{i:f_i([w])>0\}$. Orthogonality implies $S(w)\cap S(v)=\varnothing$.

\begin{lemma}[Positive extension]\label{lem:positive}
Let $\mu$ be the unitarily invariant probability measure on the unit sphere of $\C^3$. Define $H(0)=0$ and $H(w)=\norm w^2 f([w])$ for $w\ne0$. For a positive semidefinite Hermitian matrix $Q$ put
\begin{equation}\label{eq:extension}
 E(Q)=3\int H(Q^{1/2}u)\,d\mu(u)\in\R^m.
\end{equation}
Then $E$ is continuous, nonnegative coordinatewise, and satisfies
\begin{equation}\label{eq:trace-support}
 \sum_iE_i(Q)=\tr Q,\qquad
 \{i:E_i(Q)>0\}=\bigcup_{0\ne w\in\ran Q}S(w).
\end{equation}
The union is empty when $Q=0$.
\end{lemma}
\begin{proof}
Since $\norm{H(w)}_1=\norm w^2$, $H$ is continuous also at zero. The positive semidefinite square root is continuous. Joint continuity of the integrand, uniformly on the unit sphere in a compact neighborhood of each $Q$, proves continuity of $E$. Nonnegativity is immediate. Unitary invariance gives $\int uu^*\,d\mu(u)=I/3$, whence
\[
 \sum_iE_i(Q)=3\int\norm{Q^{1/2}u}^2\,d\mu(u)=\tr Q.
\]
Every nonzero $Q^{1/2}u$ belongs to $\ran Q$. Conversely, if $0\ne w\in\ran Q$ and $f_i([w])>0$, take the inverse square root on $\ran Q$ and normalize $Q^{-1/2}w$ to a unit vector $u$. Then $Q^{1/2}u$ is a positive multiple of $w$, so the $i$th integrand is positive at $u$. By continuity it is positive on a nonempty open subset of the ambient unit sphere, which has positive $\mu$-measure. This proves the support identity even for singular $Q$; no ambient positive measure is asserted for its range itself.
\end{proof}

\begin{proposition}\label{prop:lower}
Every cover of $X$ by strictly smaller-diameter subsets has at least nine members.
\end{proposition}
\begin{proof}
For such a finite cover construct $E$ as above. For $A\in\Herm_3$ let
\[
 |A|=(A^2)^{1/2},\qquad A_\pm=\frac{|A|\pm A}{2},\qquad
 F(A)=E(A_+)-E(A_-).
\]
The map $F$ is continuous and odd. The ranges of $A_+$ and $A_-$ are orthogonal, so the support identity in \eqref{eq:trace-support} and the disjointness of the $S(w)$ show that the positive coordinates of $E(A_+)$ and $E(A_-)$ are disjoint. There is no cancellation. Therefore
\[
 \norm{F(A)}_1=\tr A_++\tr A_-=\tr|A|>0\quad(A\ne0).
\]
Normalization by the Euclidean norm gives a continuous odd map from the unit sphere of $\Herm_3$, which is $S^8$, to $S^{m-1}$. The Borsuk--Ulam theorem forbids such a map if $m\le8$; equivalently, compose with an equatorial inclusion into $S^7$. The empty family cannot cover nonempty $X$. This proves the assertion.
\end{proof}
Since every subset of $X$ has diameter at most $\sqrt2$, Proposition~\ref{prop:lower} is exactly the arbitrary-subset lower assertion of Theorem~\ref{thm:main}.

\section{The Hesse configuration and its twelve equality rays}
Indices in this section lie in $\mathbb Z/3\mathbb Z$. Put $\omega=e^{2\pi i/3}$ and define unitary operators
\[
 (Uv)(j)=v(j-1),\qquad V=\operatorname{diag}(1,\omega,\omega^2).
\]
Thus $VU=\omega UV$ and $U^3=V^3=I$. We use $U,V$ for shift and clock to distinguish the shift from the projector space $X$. Let
\begin{equation}\label{eq:sic}
 s_0=\frac{(0,1,-1)}{\sqrt2},\qquad s_{a,b}=U^aV^bs_0,
 \qquad \mathcal S=\{s_{a,b}:a,b\in\mathbb Z/3\}.
\end{equation}
Explicitly, the three groups are
\[
 \frac{(0,\omega^b,-\omega^{2b})}{\sqrt2},\quad
 \frac{(-\omega^{2b},0,\omega^b)}{\sqrt2},\quad
 \frac{(\omega^b,-\omega^{2b},0)}{\sqrt2}.
\]
The twelve rays to be called \emph{holes} are the standard basis $\mathcal E=(e_0,e_1,e_2)$ and three bases $\mathcal B_c=(v_{c,0},v_{c,1},v_{c,2})$, where
\begin{equation}\label{eq:mub}
 v_{c,a}(j)=\frac{\omega^{cj^2+aj}}{\sqrt3}.
\end{equation}
For unit $s,u$ write $x_s(u)=\abs{\ip su}^2$.

\begin{lemma}\label{lem:actions}
The four displayed triples are mutually unbiased orthonormal bases. Up to unit phases,
\[
 Ve_j=e_j,\quad Ue_j=e_{j+1},\quad
 Vv_{c,a}=v_{c,a+1},\quad Uv_{c,a}=v_{c,a+c}.
\]
Both $U$ and $V$ permute the rays of $\mathcal S$ and preserve each basis by a cyclic index translation.
\end{lemma}
\begin{proof}
Within a fixed $c$, orthonormality follows from $\sum_{j=0}^2\omega^{kj}=0$ for nonzero $k$. For distinct $c,c'$, the relevant unnormalized inner product is $\sum_j\omega^{dj^2+ej}$ with $d\ne0$. In its squared modulus set $j=k+t$. For $t\ne0$, summation over $k$ vanishes because $2dt\ne0$; the $t=0$ sum is three. Cross-basis squared overlaps are thus $1/3$. The standard basis has the same overlaps with every $v_{c,a}$. All twelve rays are distinct.
The formula for $V$ follows by adding $j$ in the exponent. For $U$ the exact formula is
\[
 Uv_{c,a}=\omega^{c-a}v_{c,a-2c}=\omega^{c-a}v_{c,a+c}.
\]
The remaining assertions follow from $VU=\omega UV$.
\end{proof}

\begin{lemma}\label{lem:overlap}
Distinct members of $\mathcal S$ have squared overlap $1/4$. At every hole $h$, the nine scores $x_s(h)$ consist of three zeros and six values $1/2$.
\end{lemma}
\begin{proof}
By covariance it suffices for the first assertion to compare $s_0$ with $s_{a,b}\ne s_0$. When $a=0$, the overlap is $(\omega^b+\omega^{2b})/2=-1/2$. For $a=1,2$ exactly one coordinate product remains, of modulus $1/2$. For the second assertion, the basis actions reduce to $e_0$ and $v_{c,0}$. At $e_0$ the statement is immediate. At $v_{c,0}=(1,\omega^c,\omega^c)/\sqrt3$, the three groups give
\[
 \frac{|1-\omega^{-b}|^2}{6},\quad
 \frac{|1-\omega^{c+b}|^2}{6},\quad
 \frac{|1-\omega^{c-b}|^2}{6}.
\]
Each group has one zero and two halves.
\end{proof}

For a hole $h$ write $T_h=\{s\in\mathcal S:s\perp h\}$ and $L(h)=h^\perp$.
\begin{lemma}[Trio incidence]\label{lem:trios}
For each one of the four bases $(b_0,b_1,b_2)$, the three sets $T_{b_j}$ are disjoint triples partitioning $\mathcal S$. Every $s\in T_{b_j}$ has score $1/2$ at the other two basis vectors. Each $s\in\mathcal S$ lies on exactly one of the three lines $L(b_j)$ from each basis. If $p,q$ are the other two basis vectors and $T_r\subset\operatorname{span}(p,q)$, phases can be chosen so that
\begin{equation}\label{eq:trio}
 s=\frac{p+\overline{\gamma_s}q}{\sqrt2}\quad(s\in T_r),
 \qquad \{\gamma_s:s\in T_r\}=\gamma\{1,\omega,\omega^2\},\quad |\gamma|=1.
\end{equation}
\end{lemma}
\begin{proof}
A member orthogonal to two basis vectors would equal the third ray, giving score one there, contrary to Lemma~\ref{lem:overlap}. The three triples are therefore disjoint and exhaust the nine members. The remaining scores are halves. This also gives the one-line-per-basis assertion. Normalize $\ip ps=1/\sqrt2$ to obtain the first expression in \eqref{eq:trio}. For distinct $s,s'$ the overlap identity gives
\[
 \frac14=\frac{|1+\gamma_s\overline{\gamma_{s'}}|^2}{4},
 \qquad \Re(\gamma_s\overline{\gamma_{s'}})=-\frac12.
\]
Three unit numbers with these pairwise real parts form a rotated cube-root triple.
\end{proof}

\begin{lemma}[Exact equality locus]\label{lem:holes}
Let $M_0(u)=\max_{s\in\mathcal S}x_s(u)$. Then $M_0(u)\ge1/2$ for every unit $u$, with equality precisely at the twelve hole rays in \eqref{eq:mub} and $\mathcal E$.
\end{lemma}
\begin{proof}
Put $p_j=|u_j|^2$. For each coordinate pair $(a,b)$ the associated SIC group gives
\[
 \max_{\zeta\in\{1,\omega,\omega^2\}}|u_a-\zeta u_b|^2
 =p_a+p_b+2\sqrt{p_ap_b}\,\mu_{ab},\qquad \tfrac12\le\mu_{ab}\le1.
\]
When both coordinates are nonzero, $\mu_{ab}$ is the maximum of three negative cosines at angles separated by $2\pi/3$. The bound follows because one angle lies within $\pi/3$ of $\pi$. If a coordinate vanishes, the term containing $\mu_{ab}$ is zero. Choose $c$ with $p_c$ minimal. Then
\[
 M_0(u)\ge\frac{1-p_c+\sqrt{p_ap_b}}2\ge\frac12.
\]
If equality holds, $\sqrt{p_ap_b}=p_c$. If $p_c=0$, then $p_ap_b=0$, so $u$ is a coordinate ray. Otherwise $p_a=p_b=p_c=1/3$, and equality in every pair forces $\mu_{ab}=1/2$. The three cosine angles must then be $0,2\pi/3,4\pi/3$ up to permutation. Hence all relative coordinate phases are cube roots of unity. After choosing the first coordinate positive,
\[
 u=(1,\omega^\alpha,\omega^\beta)/\sqrt3.
\]
The map $(c,a)\mapsto(c+a,c+2a)$ is invertible modulo three, so these nine rays are exactly the $v_{c,a}$. Conversely Lemma~\ref{lem:overlap} gives equality at each of the twelve listed holes.
\end{proof}
The unperturbed maximum-overlap cell for any $s\in\mathcal S$ contains the two orthogonal poles of a basis at which its score is $1/2$. Thus its diameter is $\sqrt2$; the non-strict Hesse cover is insufficient for Theorem~\ref{thm:main}.

\section{A covariant perturbation with a uniform local gap}
In each ordered basis $(b_0,b_1,b_2)$ declare that the line $L(b_a)$ \emph{favors} the pole $b_{a+1}$. Every hole is favored by exactly one line. Lemma~\ref{lem:actions} makes this assignment compatible with both $U$ and $V$.

Put
\begin{equation}\label{eq:perturb}
 d_0=\frac{1+i}{2}(0,1,1),\qquad
 f_\varepsilon=\frac{s_0+\varepsilon d_0}{\sqrt{1+\varepsilon^2}},\qquad
 s_{a,b}(\varepsilon)=U^aV^bf_\varepsilon.
\end{equation}
We have $\norm{d_0}=1$ and $\ip{s_0}{d_0}=0$. For $s=s_{a,b}$ let $d_s=U^aV^bd_0$, so $s(\varepsilon)=(s+\varepsilon d_s)/\sqrt{1+\varepsilon^2}$. Define
\[
 \kappa(s,h)=\Re\big(\ip{d_s}h\ip hs\big),\qquad
 M_\varepsilon(u)=\max_{a,b}|\ip{s_{a,b}(\varepsilon)}u|^2.
\]
Simultaneous phase changes of $s,d_s$ and phase changes of $h$ do not change $\kappa$.

\begin{lemma}[Favored derivatives]\label{lem:derivatives}
If $L(b_a)$ contains $s$ and favors $h=b_{a+1}$, then
\[
 \kappa(s,h)\in\left\{\frac1{2\sqrt2},\frac1{2\sqrt6}\right\},
 \qquad \kappa(s,h)\ge\frac1{\sqrt{24}}.
\]
\end{lemma}
\begin{proof}
WH covariance transports the attached pair $(s,d_s)$ with a common phase and preserves the favored-pole rule. It suffices to consider the four lines through $s_0$: $L(e_0)$ and $L(v_{c,0})$, favoring $e_1$ and $v_{c,1}$ respectively. Directly,
\[
 \ip{d_0}{e_1}\ip{e_1}{s_0}=\frac{1-i}{2\sqrt2}.
\]
For the other three poles,
\[
 \ip{d_0}{v_{c,1}}=-\frac{(1-i)\omega^c}{2\sqrt3},\qquad
 \ip{v_{c,1}}{s_0}=-\frac{i\omega^{-c}}{\sqrt2},
\]
whose product is $(1+i)/(2\sqrt6)$. Taking real parts proves the claim.
\end{proof}

\begin{lemma}[Full complex neighborhood estimate]\label{lem:local}
Set
\begin{equation}\label{eq:constants}
 c=\frac1{\sqrt{12}},\qquad r=\frac c6=\frac1{12\sqrt3}.
\end{equation}
Let $h=b_{a+1}$, $h'=b_{a+2}$ and $h''=b_a$ in one of the four bases. If
\[
 u=\frac{h+z_2h'+z_3h''}{N},\qquad
 N^2=1+|z_2|^2+|z_3|^2,\qquad |z|<r,
\]
then $M_\varepsilon(u)>1/2$ for every $0<\varepsilon\le1/50$.
\end{lemma}
\begin{proof}
Let $a=|z_2|$, $b=|z_3|$. For $s$ in the favored trio $T_{h''}$, normalize the attached pair $(s,d_s)$ by a common phase so that
\[
 s=(h+\overline\gamma h')/\sqrt2.
\]
Put $\delta_0=\ip{d_s}h$, $\delta_1=\ip{d_s}{h'}$, $\delta_2=\ip{d_s}{h''}$, and $D=\delta_0+\delta_1z_2+\delta_2z_3$. These are bra coefficients, so no further conjugation occurs in $D$. Orthonormality gives $\sum|\delta_j|^2=1$ and
\[
 \ip{s(\varepsilon)}u=
 \frac{(1+\gamma z_2)/\sqrt2+\varepsilon D}{N\sqrt{1+\varepsilon^2}}.
\]
Expanding the squared modulus yields the exact identity
\begin{align}\label{eq:expansion}
 &2N^2(1+\varepsilon^2)\left(x_{s(\varepsilon)}(u)-\frac12\right)\notag\\
 &\quad=2\Re(\gamma z_2)-b^2
 +2\sqrt2\,\varepsilon\Re\big(\overline{1+\gamma z_2}D\big)
 +2\varepsilon^2|D|^2-\varepsilon^2N^2.
\end{align}
Cauchy--Schwarz gives $\Re D\ge\Re\delta_0-|z|$ and, since $|z|<r<1$,
\[
 |\overline{\gamma z_2}D|\le |z|(1+|z|)\le2|z|.
\]
Thus
\begin{equation}\label{eq:crossbound}
 \Re\big(\overline{1+\gamma z_2}D\big)\ge\Re\delta_0-3|z|.
\end{equation}
The phase normalization gives $\Re\delta_0=\sqrt2\kappa(s,h)\ge c$ for every member of the favored trio. By \eqref{eq:trio}, one member satisfies $\Re(\gamma z_2)\ge a/2$ for every complex $z_2$. Since $3|z|<c/2$ and $N^2<2$, its scaled score in \eqref{eq:expansion}, denoted $E_1$, satisfies
\begin{equation}\label{eq:F}
 E_1>a-b^2+\frac{\varepsilon}{\sqrt6}-2\varepsilon^2.
\end{equation}
For the other trio $T_{h'}$, interchange $z_2,z_3$ and use only $\Re\delta_0\ge-1$. Some member has scaled score $E_2$ satisfying
\begin{align}\label{eq:O}
 E_2&\ge b-a^2-\left(2\sqrt2+\frac1{\sqrt6}\right)\varepsilon-2\varepsilon^2\notag\\
    &\ge b-a^2-\frac{13}{4}\varepsilon-2\varepsilon^2.
\end{align}
All constants can be checked exactly:
\begin{equation}\label{eq:rationalconstants}
 r^3=\frac1{5184\sqrt3}<\frac1{1000},\qquad
 \frac1{\sqrt6}>\frac25,\qquad
 2\sqrt2+\frac1{\sqrt6}<\frac{13}{4}.
\end{equation}
For the last inequality, $\sqrt2<99/70$ and $1/\sqrt6<41/100$ give the upper bound $2267/700<13/4$; both radical inequalities follow by squaring positive quantities.

If $b^2\le a$, \eqref{eq:F} gives
\[
 E_1>\varepsilon\left(\frac25-2\varepsilon\right)
 \ge\frac9{25}\varepsilon>0.
\]
If $b^2>a$, then $a^2<b^4<r^3b<b/1000$. For $b\ge4\varepsilon$, \eqref{eq:O} gives
\[
 E_2>\varepsilon\left(\frac{3996}{1000}-\frac{13}4-\frac1{25}\right)
 =\frac{353}{500}\varepsilon>0.
\]
For $b<4\varepsilon$, \eqref{eq:F} instead gives
\[
 E_1>\varepsilon\left(\frac25-18\varepsilon\right)
 \ge\frac1{25}\varepsilon>0.
\]
These cases include zero coordinates and every complex phase, as well as the boundaries $b^2=a$, $b=4\varepsilon$ and $\varepsilon=1/50$. The positive scaling factor in \eqref{eq:expansion} proves that one of the two chosen scores exceeds $1/2$.
\end{proof}

\section{The global cover and its strict diameters}
For each hole $h$ with associated other basis vectors $h',h''$, define the open projective neighborhood
\[
 B_h=\{[u]:|\ip{h'}u|^2+|\ip{h''}u|^2<r^2|\ip hu|^2\}.
\]
This is exactly the chart region $|z|<r$ of Lemma~\ref{lem:local}. Put $K=\CP^2\setminus\bigcup_hB_h$.

\begin{theorem}\label{thm:uniform}
The set $K$ is nonempty and compact, and
\[
 \eta=\min_{[u]\in K}M_0(u)-\frac12>0.
\]
For every
\begin{equation}\label{eq:epsilon}
 0<\varepsilon\le\min\left\{\frac1{50},\frac\eta3\right\},
\end{equation}
we have $M_\varepsilon(u)>1/2$ for every $[u]\in\CP^2$.
\end{theorem}
\begin{proof}
The sets $B_h$ are open and contain their respective holes, so $K$ is compact. It is nonempty: the ray $[s_0]$ lies in no $B_h$. Indeed, membership would imply
\[
 |\ip h{s_0}|^2>\frac1{1+r^2}>\frac12,
\]
whereas Lemma~\ref{lem:overlap} gives only zero or $1/2$. The minimum defining $\eta$ is thus well-defined. Lemma~\ref{lem:holes} excludes equality throughout $K$, and continuity gives $\eta>0$.
For unit $s,s',u$, factoring the difference of squared moduli gives
\[
 \big||\ip{s'}u|^2-|\ip su|^2\big|\le2\norm{s'-s}.
\]
Moreover
\[
 \norm{s_{a,b}(\varepsilon)-s_{a,b}}^2
 =2-\frac2{\sqrt{1+\varepsilon^2}}\le\varepsilon^2.
\]
For the last inequality, $t-2+2/\sqrt{1+t}$ is zero at $t=0$ and has derivative $1-(1+t)^{-3/2}\ge0$ for $t\ge0$. Consequently $M_\varepsilon\ge M_0-2\varepsilon$. On $K$, condition~\eqref{eq:epsilon} gives
\[
 M_\varepsilon\ge\frac12+\eta-2\varepsilon\ge\frac12+\frac\eta3>\frac12.
\]
On each $B_h$ apply Lemma~\ref{lem:local}. These regions cover $\CP^2$.
\end{proof}
No numerical value of $\eta$ is claimed. Its positivity is enough to choose a positive $\varepsilon$ satisfying~\eqref{eq:epsilon}.

\begin{proof}[Proof of the upper bound in Theorem~\ref{thm:main}]
Fix such an $\varepsilon$. By continuity on compact $\CP^2$,
\begin{equation}\label{eq:delta}
 \delta=\min_{[u]\in\CP^2}\left(M_\varepsilon(u)-\frac12\right)>0.
\end{equation}
For each $(a,b)$ let
\begin{equation}\label{eq:cells}
 A_{a,b}=\{P_u\in X:|\ip{s_{a,b}(\varepsilon)}u|^2=M_\varepsilon(u)\}.
\end{equation}
These nine sets are closed and cover $X$, including all ties. Each lies in the cap where the score against its designated center is at least $1/2+\delta$. If $u\perp v$, Bessel's inequality gives
\[
 |\ip su|^2+|\ip sv|^2\le1
\]
for every unit $s$, so no $A_{a,b}$ contains an orthogonal pair. Each cell is nonempty (its center belongs to it) and compact. The continuous distance attains its maximum on its square; by \eqref{eq:projector} that maximum is strictly smaller than $\sqrt2$. Hence $b(X)\le9$, completing the proof of Theorem~\ref{thm:main}.
\end{proof}

\begin{corollary}\label{cor:balls}
Nine closed Fubini--Study balls of one common radius strictly smaller than $\pi/4$ cover $\CP^2$.
\end{corollary}
\begin{proof}
Using \eqref{eq:delta}, take the nine centers from \eqref{eq:perturb} and radius $\rho=\arccos\sqrt{1/2+\delta}<\pi/4$. The maximum-score inequality means that each ray belongs to at least one such ball.
\end{proof}

\section{Scope of the construction and comparison}\label{sec:comparison}
The Hesse SIC itself has $\alpha=1/2$ by Lemma~\ref{lem:holes}. Its unperturbed maximum-overlap cells contain orthogonal pairs; neither its optimal separation nor its incidence data makes their diameters strictly smaller. The perturbation deliberately leaves the SIC locus. If $p,q$ are the two nonzero squared coordinate moduli of $f_\varepsilon$, then
\[
 p=\frac12+\frac{\varepsilon}{\sqrt2(1+\varepsilon^2)},\qquad
 q=\frac12-\frac{\varepsilon}{\sqrt2(1+\varepsilon^2)},\qquad
 pq=\frac14-\frac{\varepsilon^2}{2(1+\varepsilon^2)^2}.
\]
Within a coordinate-zero group, distinct orbit rays have squared overlap $1-3pq>1/4$; between groups it is $pq<1/4$. Thus the positive-$\varepsilon$ configuration is not a unitary relabeling of a SIC or of a configuration having twenty-seven overlaps exactly equal to $1/4$.

Hughston--Salamon classify qutrit SICs without assuming covariance~\cite{HS}. In their midpoint representation all nonzero coordinates of a SIC ray have equal modulus. A coordinate ray therefore has maximum score $1/2$, so that classified family does not by itself give the strict cap cover here. Their explicit non-SIC examples require a separate comparison: Appendix~\ref{app:nonsic} gives exact witnesses below $1/2$ for the three-basis WH example and for the particular quartic example in their Section~10. This does not exclude every choice in their general polynomial family or every non-SIC WH orbit.

A different useful distinction is provided by Mondal--Heath~\cite[Theorem~1 and Appendix~A]{MH}. Their covering-codebook criterion identifies projection covering radius strictly less than one with covering by complementary subspaces. For rays, this gives $\alpha>0$; the requirement here is $\alpha>1/2$. Compactness supplies a positive margin from pointwise positivity, but does not prescribe that margin to exceed $1/2$. Their coordinate antenna codebook has $\alpha=1/3$. Likewise, for any codebook whose normalized vectors have all three coordinate moduli $1/\sqrt3$, every coordinate ray has score $1/3$ against every center. A packing assertion about such a codebook cannot give the present cover.

The expected-distortion quantifier in complex Grassmann quantization, for example $\mathbb E_Q\min_{s\in C}d_c(Q,s)^2$ in Dai--Liu--Rider~\cite{DLR}, differs from the maximum over every target ray for a fixed nine-center codebook. No inference from that expectation to \eqref{eq:alpha} is made. The comparison here is confined to the stated primary constructions and operations. It does not assert that every older finite covering table or codebook family has been exhausted, or that our covering radius is optimal.

\section{Verification scope and reproducibility}\label{sec:verification}
The proof above is a human mathematical proof. Its global gap $\eta$ is existential, and no sampled minimum or numerical choice of $\varepsilon$ is used. The accompanying reproducibility package fixes the source identities and records the exact scalar checks, manuscript build and page inspection.

A separate Lean development in the namespace \texttt{BorsukLocal} checks the scalar three-case positivity argument and exact complex norm-square expansions used in Lemma~\ref{lem:local}. In particular, for $a,b\ge0$, $b\le1/10$ and $0<\varepsilon\le1/50$, its \texttt{two\_trio\_margin} theorem proves
\[
 \frac{\varepsilon}{25}<\max\left\{
 a-b^2+\frac25\varepsilon-2\varepsilon^2,
 b-a^2-\frac{13}4\varepsilon-2\varepsilon^2\right\}.
\]
The development also checks the displayed radical bounds and scalar conversions from scaled score positivity. It does not formalize the selection of the two trios, their geometric phase normalizations, the Cauchy--Schwarz estimates producing \eqref{eq:F}--\eqref{eq:O}, the Hesse/WH incidence geometry, compactness, the Hermitian integral extension, or Borsuk--Ulam. Theorem~\ref{thm:main} is therefore not claimed to be fully formalized. Exact coverage and verification receipts for this limited support are recorded separately in the accompanying provenance file.

\appendix
\section{Two explicit non-SIC configurations}\label{app:nonsic}
This appendix clarifies specific potential antecedents. A witness with all nine scores at most $1/2$ excludes the sufficient cap-cover property of that codebook, not arbitrary alternative partitions of $\CP^2$.

\subsection{A WH orbit containing orthogonal rays}
\begin{proposition}\label{prop:whorth}
If a nine-ray qutrit Weyl--Heisenberg orbit contains an orthogonal pair of distinct rays, its value of $\alpha$ is at most $1/3$.
\end{proposition}
\begin{proof}
Write the orbit as $\{U^aV^bf\}$. An orthogonal pair implies $\ip f{Df}=0$ for a nonidentity displacement $D$, up to an irrelevant scalar phase. Its three eigenvalues, after a common phase change, are $1,\omega,\omega^2$, all simple. If $p_j$ are the squared moduli of $f$ in that eigenbasis, then
\[
 p_0+p_1+p_2=1,\qquad p_0+\omega p_1+\omega^2p_2=0.
\]
Taking real and imaginary parts gives $p_0=p_1=p_2=1/3$. Conjugation of $D$ by any WH displacement multiplies it by a scalar, so the same reasoning applies to every orbit vector. Each $D$-eigenvector consequently has squared overlap $1/3$ with all nine rays.
\end{proof}
Hughston--Salamon's Example~6.4~\cite{HS} is a WH orbit forming three pairwise mutually unbiased orthonormal bases. Its twenty-seven nonzero pair scores are $1/3$ and its other nine pairs are orthogonal. Proposition~\ref{prop:whorth} excludes the strict cap-cover property for that example without assuming that an arbitrary MUB triple extends to a fourth basis.

\subsection{The explicit quartic example}
In the coordinates of~\cite[(6.3)]{HS}, put
\[
 z(\sigma,x)=\frac1{\sqrt{1+x^2}}
 \left(\sqrt{\frac23}e^{i\sigma},
       -\frac1{\sqrt6}-\frac{x}{\sqrt2},
       -\frac1{\sqrt6}+\frac{x}{\sqrt2}\right).
\]
Example~10.2 has two fixed rays
\[
 z_1=(0,1,-\omega)/\sqrt2,\qquad z_2=(0,1,-\omega^2)/\sqrt2,
\]
and seven others $z(0,t)$ and $z(\pm\sigma_j,x_j)$ for $j=4,5,6$. Here $t=x_3<x_4<x_5<x_6$ are the roots of
\begin{equation}\label{eq:quartic}
 27x^4-66x^2+8\sqrt{26-2\sqrt{97}}\,x+3=0,
\end{equation}
and $\sigma_j>0$ makes the squared overlap with $z(0,t)$ equal to $1/4$. Expanding that overlap gives
\begin{equation}\label{eq:cos}
 \cos\sigma_j=
 \frac{9(1+t^2)(1+x_j^2)-16-4(1+3tx_j)^2}{16(1+3tx_j)}.
\end{equation}
Take $u=(1,0,1)/\sqrt2$. Its scores against $z_1,z_2$ are $1/4$, while
\begin{equation}\label{eq:witnessscore}
 |\ip u{z(\sigma,x)}|^2=
 \frac{4+(\sqrt3x-1)^2+4(\sqrt3x-1)\cos\sigma}{12(1+x^2)}.
\end{equation}
The sign of $\sigma$ does not matter.

For completeness, the rational intervals used to certify the bound are
\[
\begin{array}{c|c}
 t&[-1.687100,-1.687097]\\
 x_4&[-0.109890,-0.109887]\\
 x_5&[0.442461,0.442464]\\
 x_6&[1.354524,1.354527].
\end{array}
\]
Every decimal endpoint in this appendix denotes its exact terminating rational value. The radical bounds
\begin{align*}
 9.8488578&<\sqrt{97}<9.8488579,\\
 2.5104350&<\sqrt{26-2\sqrt{97}}<2.5104352,\\
 1.7320508&<\sqrt3<1.7320509
\end{align*}
are verified by squaring positive quantities. Rational interval substitution makes the polynomial in \eqref{eq:quartic} have opposite signs at the two ends of each of the four disjoint intervals. They therefore contain all four roots, one each. Substitution into \eqref{eq:cos} and \eqref{eq:witnessscore}, with denominator intervals avoiding zero, bounds the four different scores respectively in
\[
 (0.080,0.081),\quad(0.298,0.300),\quad
 (0.345,0.347),\quad(0.263,0.265).
\]
These inclusions can be checked with rational arithmetic; an accompanying standard-library Fraction script records the full exact fraction endpoints. In particular all nine scores are strictly less than $7/20$, and this explicit codebook satisfies $\alpha<7/20<1/2$. No numerical minimization or asserted optimality of the witness enters this conclusion.

The printed parametrization in~\cite[(7.1)]{HS} multiplies the last coordinate by $\omega^2$. If that convention is used literally, multiply the last coordinate of $u$ by $\omega^2$ as well. The seven scores are unchanged and the scores against $z_1,z_2$ remain $1/4$, since the middle coordinate of $u$ is zero. The exclusion holds under either convention. It concerns the explicit Example~10.2, not every parameter choice in the paper's general case-(iv) family.

\section*{Rights and provenance}
The author has no institutional affiliation. Original prose is licensed under Creative Commons Attribution 4.0; original accompanying code is licensed under MIT. Cited third-party materials retain their own rights. The lower-bound method is credited in Section~1 and reproved in Section~2. The source version used is fixed by the repository commit in reference~\cite{OpenAI}.

\begin{thebibliography}{9}
\raggedright
\bibitem{OpenAI}
OpenAI, \emph{A nine-dimensional counterexample to Borsuk's covering assertion}, September 23, 2026. Repository \texttt{openai/math}, revision \texttt{fd4aeeb2ee4fc729c18d98444fed42fd0529eeeb}, especially the geometry, extension and topology sections.\newline
\href{https://github.com/openai/math/tree/fd4aeeb2ee4fc729c18d98444fed42fd0529eeeb/preprints/A-nine-dimensional-counterexample-to-Borsuks-covering-assertion-September-23-2026}{Pinned source and companion files}.
\bibitem{Dang}
H. B. Dang, K. Blanchfield, I. Bengtsson and D. M. Appleby,
\emph{Linear dependencies in Weyl--Heisenberg orbits},
\href{https://arxiv.org/abs/1211.0215v2}{arXiv:1211.0215v2}.
\bibitem{HS}
L. Hughston and S. Salamon,
\emph{Surveying points in the complex projective plane},
\href{https://arxiv.org/abs/1410.5862v2}{arXiv:1410.5862v2} (2015).
\bibitem{MH}
B. Mondal and R. W. Heath Jr.,
\emph{A diversity guarantee and SNR performance for unitary limited feedback MIMO systems},
EURASIP Journal on Advances in Signal Processing (2008), Article 594928, 15 pp.
\href{https://doi.org/10.1155/2008/594928}{doi:10.1155/2008/594928}.
\bibitem{DLR}
W. Dai, Y. Liu and B. Rider,
\emph{Quantization bounds on Grassmann manifolds and applications to MIMO communications},
\href{https://arxiv.org/abs/cs/0603039}{arXiv:cs/0603039}.
\end{thebibliography}
\end{document}
